\section{Analisis vertical}

El analisis vertical expresa cada partida como porcentaje de una base comun dentro del mismo estado financiero. Permite estudiar la estructura economica o financiera de la empresa.

La formula es:

\[
\text{Participacion porcentual} =
\frac{\text{Partida}}{\text{Base}} \times 100
\]

En el Estado de Situacion Financiera, la base usual es el total activo. En el Estado de Resultados, la base usual son las ventas netas.

Este metodo responde a la pregunta: \textbf{que peso tiene una partida dentro del total?}

\subsection{Ejemplo basico de analisis vertical}

Enunciado: una empresa desea conocer que proporcion de sus ventas se destina al costo de ventas.

Datos:

\begin{itemize}
  \item Ventas netas: S/ 100,000.00.
  \item Costo de ventas: S/ 60,000.00.
\end{itemize}

Calculo:

\[
\frac{S/\,60,000.00}{S/\,100,000.00} \times 100 = 60.00\%
\]

Interpretacion: de cada S/ 100.00 vendidos, S/ 60.00 se destinan al costo de ventas y S/ 40.00 quedan como margen bruto antes de gastos.

\subsection{Ejemplo aplicado de analisis vertical}

Enunciado: una empresa tecnologica desea evaluar que peso tienen sus costos de servidores y licencias dentro de sus ventas.

Datos:

\begin{itemize}
  \item Ventas: S/ 200,000.00.
  \item Servidores y licencias: S/ 70,000.00.
\end{itemize}

Calculo:

\[
\frac{S/\,70,000.00}{S/\,200,000.00} \times 100 = 35.00\%
\]

Interpretacion y decision: el 35.00\% de las ventas se destina a infraestructura tecnologica. Si este porcentaje aumenta, la empresa debe revisar arquitectura, proveedores cloud, licencias y escalabilidad.
